\chapter{Conclusion and Future Work}

\section{Conclusion}

This project set out to build an EEG system around a four-channel dry-electrode headset
that can tell left-hand from right-hand motor imagery and, from the same recording,
report the mental state of the person as a cognitive score.

The motor imagery analysis was built and checked in three stages. First, the ERD/ERS
analysis of the published stroke patient dataset was reproduced, which confirmed that the
processing chain gives the same result as the original paper. Second, the same method was
applied to all nine healthy subjects of the BCI Competition IV 2b dataset, which matches
the headset in sampling rate and electrode positions. Here the expected pattern appeared
clearly: the 8--30 Hz power dropped more over the side of the head opposite to the
imagined hand in 16 of 18 cases, and subject B04 showed the textbook result, with C4
falling to about $-50$ per cent during left-hand imagery and C3 to about $-60$ per cent
during right-hand imagery. Third, the analysis was applied to left and right motor imagery
recorded with the GELO headset.

The most useful single finding is the C3/C4 power ratio. On the 2b data it was higher for
the left hand than for the right hand in all nine subjects, and as a single number per
trial it identified the hand on unseen days with 73.5 per cent accuracy, better than chance
for every subject and 96.6 per cent for B04. It works better than the power or the ERD of
either channel on its own because it cancels changes that affect both channels together
and needs no rest baseline. The best classifiers reached 80.2 per cent when trained on the
same person and 74.7 per cent on a person never seen in training, using filter bank CSP or
simple log band power with an SVM and Euclidean alignment. With only two or three channels
the choice of classifier made very little difference, while Euclidean alignment helped
every model.

On the data recorded with the GELO headset, the same kind of C3/C4 difference is visible in
the trials examined, with a power ratio of 9.52 for a left-hand trial and 0.86 for a
right-hand trial. However, classifiers trained on either public dataset stayed close to
chance on the recorded data, and across all 72 usable trials the left and right difference
was much weaker than in the 2b data. Since the same classifier works well on new people in
2b, the main limit is the recorded data and the way it was collected, and not the
classifier or the number of channels.

The cognitive assessment part works end to end. The electrode contact settles within a few
minutes, the usable measurement band of the headset was found to be about 4 to 30 Hz, and
attention, cognitive load, relaxation, stress and fatigue are computed live and compared
against each person's own calibration, so that changes in overall signal level do not
move the scores.

\section{Future Work}

The results point clearly to what should be done next, in order of expected benefit.

\begin{itemize}
    \item \textbf{A better recording protocol.} New GELO recordings should follow the BCI
          Competition IV 2b design: left and right trials mixed randomly within one session,
          cues and markers placed by software at fixed times, a calm fixation period before
          each cue to act as a true baseline, equal numbers of trials for each hand, and
          several sessions on different days.
    \item \textbf{Per-user calibration.} Training on the same person's earlier sessions was
          worth about 5 percentage points over a model trained on other people, and the
          C3/C4 threshold differs between people. A short calibration run of cued trials at
          the start of use is a cheap way to gain this.
    \item \textbf{Testing the 2b model again on new GELO data.} Once data recorded with the
          new protocol is available, the 2b model with Euclidean alignment can be applied to
          it directly, with the 74.7 per cent reached on new 2b subjects as the reference.
    \item \textbf{Real-time operation on the Raspberry Pi 5.} The C3/C4 ratio and the log
          band power SVM are very light to compute, so both the motor imagery decision and
          the cognitive score can run live on the embedded platform described in
          Chapter 4.
    \item \textbf{Deep learning and imagined speech.} Models such as convolutional
          transformers need far more data than is available now. With a larger and cleaner
          dataset recorded on the headset, they become worth trying, both for motor imagery
          and for the longer-term aim of decoding imagined words.
\end{itemize}
